\changecolours{ScoutGreen}
\section[Revisão e Objetivos]{Ponte Pedagógica e Objetivos de Aprendizagem}

% ------------------------------------------------------------------------------
\twocolframe[title={Do Rascunho em Papel (Aula 07) ao Artefato CASE (Aula 08)},leftcol=7.0cm,rightcol=7.0cm]{%
  \alert{O Que Consolidamos na Aula 07 (Sábado)?}\par
  \itemseps[topsepi=0.5mm,itemsepi=1mm,labelsep=2mm,leftmargini=3mm]
  \begin{itemize}
    \item \textbf{Escopo da AVP1:} Narrativa completa das regras de negócio do projeto semestral.
    \item \textbf{Catálogo Completo:} Entidades regulares, atributos de Chen, relacionamentos ($1:1, 1:N, N:M$, recursivo), entidade fraca e hierarquia EER.
    \item \textbf{Esboço Manual:} Desenho preliminar à mão na folha do caderno de estudo dirigido.
  \end{itemize}
}{%
  \alert{A Missão Prática da Aula 08}\par
  \itemseps[topsepi=0.5mm,itemsepi=1mm,labelsep=2mm,leftmargini=3mm]
  \begin{itemize}
    \item \textbf{Digitalização Profissional:} Transformar o esboço manual em um Diagrama Entidade-Relacionamento Estendido (DER/EER) impecável usando \textbf{Ferramentas CASE}.
    \item \textbf{Bilinguismo Visual:} Dominar tanto o padrão internacional de \textbf{Peter Chen (Draw.io)} quanto o padrão brasileiro de \textbf{Heuser (brModelo)}.
    \item \textbf{Dicionário de Dados:} Documentar em tabelas descritivas as regras semânticas do modelo.
  \end{itemize}
}

% ------------------------------------------------------------------------------
\twocolframe[title={Objetivos de Aprendizagem da Aula 08},leftcol=7.0cm,rightcol=7.0cm]{%
  \alert{Competências Conceituais e Notacionais}\par
  \itemseps[topsepi=0.5mm,itemsepi=1mm,labelsep=2mm,leftmargini=3mm]
  \begin{itemize}
    \item \textbf{Ferramentas CASE:} Diferenciar editores gráficos puros de ferramentas baseadas em metamodelo de dados.
    \item \textbf{Tradução Chen $\leftrightarrow$ Heuser:} Converter elipses e linhas duplas para pinos identificadores e pares $(min, max)$.
    \item \textbf{Especialização no brModelo:} Aplicar a dupla de restrições $(t/p, e/c)$ equivalente a $(\text{Total/Parcial}, d/o)$.
  \end{itemize}
}{%
  \alert{Competências Práticas de Modelagem}\par
  \itemseps[topsepi=0.5mm,itemsepi=1mm,labelsep=2mm,leftmargini=3mm]
  \begin{itemize}
    \item \textbf{Operação no brModelo e Draw.io:} Construir diagramas limpos, alinhados em grade e sem cruzamento de arestas.
    \item \textbf{Prevenção de Antipadrões:} Evitar chaves estrangeiras no nível conceitual e entidades desconectadas.
    \item \textbf{Modelagem Descritiva e Textual:} Elaborar o Dicionário de Dados Conceitual e a Matriz de Rastreabilidade.
  \end{itemize}
}
